% Source: https://artofproblemsolving.com/community/c3h1282775_geometry_power_of_a_point
% Stable identity: aops-topic-1282775-q01
Let $M$ be the midpoint of $O_1O_2$ and $d=O_1O_2$. By the median-length formula (or by vectors with origin $M$),
\[
XO_1^2+XO_2^2=2XM^2+\dfrac{d^2}{2}.
\]

\begin{asy}
usepackage("amsmath");
size(5.5cm);
defaultpen(fontsize(10pt));
pair O1=(0,0),O2=(4,0),M=(2,0); draw(circle(O1,1.2),blue); draw(circle(O2,0.8),blue); draw(circle(M,1.8),red+dashed);
dot("$O_1$",O1,S); dot("$O_2$",O2,S); dot("$M$",M,S); pair X=M+1.8*dir(70); dot("$X$",X,N); draw(X--O1,gray); draw(X--O2,gray);
label("\mbox{}",(min(currentpicture,user=true)+max(currentpicture,user=true))/2,black); // keeps the text after the figure black when LaTeX typesets the labels inline
\end{asy}

So the condition $XO_1^2-r_1^2+XO_2^2-r_2^2=k$ becomes
\[
XM^2=\dfrac{k+r_1^2+r_2^2}{2}-\dfrac{d^2}{4}=:\rho .
\]
Hence the locus is
\begin{itemize}
\item the circle with centre $M$ (the midpoint of $O_1O_2$) and radius $\sqrt\rho$ if $\rho>0$;
\item the single point $M$ if $\rho=0$;
\item empty if $\rho<0$. $\blacksquare$
\end{itemize}
